\documentclass[12pt,oneside]{article}

%% ── Packages ──────────────────────────────────────────────────
\usepackage[T1]{fontenc}
\usepackage[utf8]{inputenc}
\usepackage{mathpazo}          % Palatino text + math
\usepackage[scaled=0.92]{helvet}
\usepackage[margin=1.1in,top=1.2in,bottom=1.25in]{geometry}
\usepackage{microtype}
\usepackage{xcolor}
\usepackage{titlesec}
\usepackage{enumitem}
\usepackage{hyperref}
\usepackage{booktabs}
\usepackage{array}
\usepackage{tabularx}
\usepackage{longtable}
\usepackage{fancyhdr}
\usepackage{amsmath,amssymb}
\usepackage{braket}
\newcommand{\proj}[1]{|#1\rangle\langle#1|}



\usepackage{tcolorbox}
\usepackage{parskip}
\usepackage{multicol}
\usepackage{setspace}
\usepackage{ragged2e}

%% ── Colour Palette ───────────────────────────────────────────
\definecolor{primary}{HTML}{1A3A5C}
\definecolor{accent}{HTML}{2672B0}
\definecolor{highlight}{HTML}{EBF4FF}
\definecolor{partbg}{HTML}{1A3A5C}
\definecolor{lightgray}{HTML}{F4F5F6}
\definecolor{rulegray}{HTML}{BFC9CA}
\definecolor{warmgold}{HTML}{9A7D0A}
\definecolor{darktext}{HTML}{1C2833}

%% ── Hyperref ──────────────────────────────────────────────────
\hypersetup{
  colorlinks=true,
  linkcolor=accent,
  urlcolor=accent,
  pdftitle={Quantum Error Correction — Comprehensive Graduate Course},
  bookmarksnumbered=true,
  bookmarksopen=false
}

%% ── Page style ────────────────────────────────────────────────
\pagestyle{fancy}
\fancyhf{}
\fancyhead[L]{\small\color{primary}\textit{Quantum Error Correction --- Comprehensive Graduate Course}}
\fancyhead[R]{\small\color{primary}\thepage}
\renewcommand{\headrulewidth}{0.4pt}
\renewcommand{\headrule}{\color{rulegray}\hrule width\headwidth height\headrulewidth}
\setlength{\headheight}{15pt}

%% ── tcolorbox styles ─────────────────────────────────────────
\tcbuselibrary{skins,breakable}

%% Part header box
\newcommand{\parthead}[3]{%
  \vspace{16pt}
  \begin{tcolorbox}[
    enhanced,colback=partbg,colframe=partbg,
    arc=3pt,left=10pt,right=10pt,top=7pt,bottom=7pt,boxrule=0pt
  ]
    \color{white}\large\bfseries
    Part #1\enspace\textbar\enspace #2
    \hfill\normalfont\normalsize\color{white!80!partbg}Lectures #3
  \end{tcolorbox}
  \vspace{2pt}
}

%% Lecture title box
\newcommand{\lecturetitle}[3]{%
  \vspace{10pt}
  \begin{tcolorbox}[
    enhanced,breakable,
    colback=highlight,colframe=accent,
    arc=2pt,left=10pt,right=8pt,top=7pt,bottom=7pt,boxrule=0.7pt
  ]
    \textbf{\color{accent}Lecture #1.\enspace}%
    {\color{primary}\bfseries #2}%
    \hfill{\small\color{warmgold}\textit{#3}}
  \end{tcolorbox}
  \vspace{-2pt}
}

%% Bullet list
\newcommand{\topics}[1]{%
  \begin{itemize}[leftmargin=1.5em,itemsep=1.5pt,topsep=3pt,parsep=0pt]
    #1
  \end{itemize}\vspace{2pt}
}
\newcommand{\T}[1]{\item #1}

%% Outcome box
\newcommand{\keybox}[1]{%
  \begin{tcolorbox}[
    enhanced,colback=lightgray,colframe=rulegray,
    arc=2pt,left=8pt,right=8pt,top=5pt,bottom=5pt,boxrule=0.4pt,
    fontupper=\small
  ]
    \textbf{\color{primary}Learning Outcomes:\enspace}#1
  \end{tcolorbox}\vspace{4pt}
}

%% Section styling (used for appendix headings)
\titleformat{\section}[block]
  {\large\bfseries\color{primary}}{}{0pt}{}
  [\vspace{-2pt}\color{accent}\rule{\textwidth}{1pt}]
\titlespacing{\section}{0pt}{18pt}{8pt}

%% ─────────────────────────────────────────────────────────────
\begin{document}

%% ═══════════════════════════════════════════════════════
%%   TITLE PAGE
%% ═══════════════════════════════════════════════════════
\thispagestyle{empty}
\begin{center}
  \vspace*{1.4cm}
  {\color{primary}\rule{\textwidth}{2pt}}
  \vspace{0.4cm}

  {\fontsize{26}{32}\selectfont\bfseries\color{primary}
   Quantum Error Correction}

  \vspace{0.2cm}
  {\large\color{accent}
   A Comprehensive Graduate Course in Theory, Algorithms, and Modern Practice}

  \vspace{0.35cm}
  {\color{primary}\rule{\textwidth}{2pt}}

  \vspace{1.1cm}

  \begin{tabular}{rl}
    \color{primary}\textbf{Format:}       & 45 Lectures $\times$ 80 minutes \\[3pt]
    \color{primary}\textbf{Level:}        & Graduate / Advanced Undergraduate \\[3pt]
    \color{primary}\textbf{Credits:}      & 4 \\[3pt]
    \color{primary}\textbf{Prerequisites:} &
       Dirac notation, density matrices (introductory), \\
    & quantum gates and circuits, linear algebra, \\
    & elementary probability and combinatorics \\[3pt]
    \color{primary}\textbf{Recommended:}  &
       Classical information theory, group theory (basics) \\
  \end{tabular}

  \vspace{1.1cm}

  \begin{tcolorbox}[
    enhanced,colback=highlight,colframe=accent,
    arc=3pt,left=12pt,right=12pt,top=10pt,bottom=10pt,
    boxrule=0.8pt,width=0.94\textwidth
  ]
    \textbf{\color{primary}Course Description.}\enspace
    This course is a rigorous, self-contained treatment of quantum error correction (QEC)
    from mathematical foundations to the state of the art.
    We develop every major theoretical framework---operator-sum representations, the
    Knill--Laflamme conditions, stabilizer formalism, topological codes, and quantum
    LDPC codes---and pair each with the algorithms and software tools needed to use
    them in practice.
    Students will leave able to design codes, prove correctness conditions, implement
    and benchmark decoders, analyse fault-tolerant architectures, and read current
    research literature with confidence.
    The course expands the standard 30-lecture treatment to 45 lectures so that
    \emph{nothing significant is omitted}: every major algorithm is derived, not just
    stated, and each code family is studied in depth.
  \end{tcolorbox}

  \vspace{1.0cm}

  \begin{tcolorbox}[
    enhanced,colback=lightgray,colframe=rulegray,
    arc=3pt,left=12pt,right=12pt,top=9pt,bottom=9pt,
    boxrule=0.4pt,width=0.94\textwidth
  ]
    \textbf{\color{primary}Course Structure}\\[6pt]
    \begin{tabularx}{\linewidth}{@{}lX r@{}}
      \color{accent}\textbf{Part I}   &
        Mathematical and Physical Foundations & Lec.\ 1--7 \\[1pt]
      \color{accent}\textbf{Part II}  &
        Elementary Quantum Codes              & Lec.\ 8--14 \\[1pt]
      \color{accent}\textbf{Part III} &
        Stabilizer Formalism                  & Lec.\ 15--22 \\[1pt]
      \color{accent}\textbf{Part IV}  &
        Fault-Tolerant Quantum Computing      & Lec.\ 23--30 \\[1pt]
      \color{accent}\textbf{Part V}   &
        Topological and Surface Codes         & Lec.\ 31--37 \\[1pt]
      \color{accent}\textbf{Part VI}  &
        Quantum LDPC Codes                    & Lec.\ 38--41 \\[1pt]
      \color{accent}\textbf{Part VII} &
        Bosonic Codes and Hardware            & Lec.\ 42--44 \\[1pt]
      \color{accent}\textbf{Part VIII}&
        Research Frontiers                    & Lec.\ 45 \\
    \end{tabularx}
  \end{tcolorbox}
\end{center}

\newpage

%% Table of Contents
{%
  \hypersetup{linkcolor=primary}
  \setlength{\parskip}{1pt}
  \tableofcontents
}

\newpage

%% ═══════════════════════════════════════════════════════════════
%%   PART I — MATHEMATICAL AND PHYSICAL FOUNDATIONS
%% ═══════════════════════════════════════════════════════════════

\parthead{I}{Mathematical and Physical Foundations}{1--7}
\addcontentsline{toc}{section}{Part I: Mathematical and Physical Foundations (Lec.\ 1--7)}

%%──────────────────────────────────────────────────────
\lecturetitle{1}{The Case for Quantum Error Correction}{Motivation}
\addcontentsline{toc}{subsection}{Lec.\ 1: The Case for Quantum Error Correction}

\topics{
  \T{Fragility of quantum information: superposition, entanglement, decoherence}
  \T{No-cloning theorem: proof and implications for redundancy}
  \T{Measurement collapse: why classical repetition fails directly}
  \T{Key insight: errors form a \emph{discrete algebra}---correcting a Pauli basis suffices for all errors (error digitisation)}
  \T{Historical milestones: Shor (1995), Steane (1996), Calderbank--Shor (1996), Kitaev (1997, 2003)}
  \T{High-level encode--detect--correct cycle}
  \T{Threshold theorem: informal statement---below a critical error rate, arbitrarily long computation is possible}
  \T{Roadmap of the course: every major concept previewed}
}

\keybox{Articulate the no-cloning obstruction; explain error digitisation; describe the QEC cycle at an intuitive level; state the threshold theorem informally.}

%%──────────────────────────────────────────────────────
\lecturetitle{2}{Quantum States and Density Operators}{Formalism I}
\addcontentsline{toc}{subsection}{Lec.\ 2: Quantum States and Density Operators}

\topics{
  \T{State vectors, superposition, and the Hilbert space $\mathcal{H}$}
  \T{Density operators $\rho$: definition, positivity, unit trace, purity}
  \T{Pure vs.\ mixed states; ensemble interpretation}
  \T{Bloch sphere: pure-state parameterisation, mixed states inside the sphere}
  \T{Tensor products: multi-qubit Hilbert spaces}
  \T{Partial trace: reduced density matrices, marginals}
  \T{Schmidt decomposition and entanglement quantification}
  \T{Purification: every mixed state is the partial trace of a pure state}
}

\keybox{Compute density matrices, reduced states, Schmidt rank, and purifications.}

%%──────────────────────────────────────────────────────
\lecturetitle{3}{Quantum Operations: CPTP Maps and Kraus Operators}{Formalism II}
\addcontentsline{toc}{subsection}{Lec.\ 3: Quantum Operations --- CPTP Maps and Kraus Operators}

\topics{
  \T{Completely positive (CP) maps: definition via positivity on extensions}
  \T{Choi--Jamio\l{}kowski isomorphism: CP maps $\leftrightarrow$ positive semidefinite Choi matrices}
  \T{Operator-sum (Kraus) representation: $\mathcal{E}(\rho)=\sum_k E_k \rho E_k^\dagger$; trace-preservation condition}
  \T{Unitary freedom in Kraus operators: two sets represent the same channel iff related by a unitary}
  \T{Stinespring dilation theorem: every CPTP map arises from unitary evolution on an extended system}
  \T{Unitary, projective, and general measurement as special cases}
  \T{Completely positive trace non-increasing maps (quantum instruments)}
}

\keybox{Derive Kraus representations; apply the Choi isomorphism; state and use the Stinespring dilation theorem.}

%%──────────────────────────────────────────────────────
\lecturetitle{4}{Standard Noise Channels in Quantum Computing}{Channels}
\addcontentsline{toc}{subsection}{Lec.\ 4: Standard Noise Channels in Quantum Computing}

\topics{
  \T{Bit-flip channel $\mathcal{E}_{BF}$: Kraus operators, action on Bloch sphere}
  \T{Phase-flip (dephasing) channel: pure dephasing $T_\phi$, relation to $T_2$}
  \T{Bit-phase-flip channel}
  \T{Depolarizing channel: isotropic Pauli noise, Bloch sphere contraction; parameter $p$}
  \T{Amplitude damping: spontaneous emission, $T_1$ relaxation; thermal amplitude damping}
  \T{Generalized amplitude damping (finite temperature)}
  \T{Erasure channel: known-location qubit loss}
  \T{Pauli channel and its importance for CSS codes and stabilizer codes}
  \T{Leakage errors: transitions outside the computational subspace}
  \T{Crosstalk: correlated multi-qubit noise (qualitative)}
  \T{Lindblad master equation: Markovian evolution and jump operators (overview)}
}

\keybox{Write Kraus operators and Bloch-sphere actions for all standard channels; connect $T_1, T_2$ to channel parameters; identify Pauli channel structure.}

%%──────────────────────────────────────────────────────
\lecturetitle{5}{Distance Measures for Quantum States and Channels}{Metrics}
\addcontentsline{toc}{subsection}{Lec.\ 5: Distance Measures for Quantum States and Channels}

\topics{
  \T{Fidelity $F(\rho,\sigma)$: Uhlmann formula, properties, pure-state case}
  \T{Bures distance, quantum angle}
  \T{Trace distance $T(\rho,\sigma)=\tfrac12\|\rho-\sigma\|_1$: operational meaning (distinguishability bound)}
  \T{Quantum relative entropy $S(\rho\|\sigma)$: definition, monotonicity under CPTP maps}
  \T{Average gate fidelity $\bar{F}$: definition and its relation to the depolarizing rate $p$}
  \T{Diamond norm $\|\mathcal{E}\|_\diamond$: definition, operational meaning (best distinguishing strategy)}
  \T{Computing diamond norm via semidefinite programming (SDP)}
  \T{Diamond norm vs.\ average gate fidelity: when they differ and which to use}
  \T{Gentle measurement lemma}
}

\keybox{Compute fidelity, trace distance, and average gate fidelity; formulate diamond norm as an SDP; explain the operational interpretation of each measure.}

%%──────────────────────────────────────────────────────
\lecturetitle{6}{Characterising Quantum Hardware Errors}{Hardware}
\addcontentsline{toc}{subsection}{Lec.\ 6: Characterising Quantum Hardware Errors}

\topics{
  \T{State preparation and measurement (SPAM) errors: modelling and mitigation}
  \T{Randomised benchmarking (RB): protocol, exponential decay, extracting $r_{EPC}$}
  \T{Interleaved RB: gate-specific error rates}
  \T{Simultaneous RB: crosstalk detection}
  \T{Quantum process tomography (QPT): protocol, $\chi$-matrix, tomographic complexity}
  \T{Gate set tomography (GST): self-consistent, SPAM-agnostic characterisation}
  \T{Pauli noise learning: sparse Pauli error models and efficient estimation}
  \T{Cycle benchmarking and Pauli error rates}
  \T{Case studies: superconducting transmons ($T_1\sim100\,\mu$s, $T_2\sim100\,\mu$s, gate time $\sim20$\,ns), trapped ions, neutral atoms}
  \T{Connecting characterised error models to code and decoder design}
}

\keybox{Design an RB or GST experiment; interpret its output; match a hardware error model to an appropriate code.}

%%──────────────────────────────────────────────────────
\lecturetitle{7}{Classical Coding Theory: Essential Background}{Classical ECC}
\addcontentsline{toc}{subsection}{Lec.\ 7: Classical Coding Theory --- Essential Background}

\topics{
  \T{Linear codes over $\mathbb{F}_2$: generator matrix $G$, parity-check matrix $H$}
  \T{Systematic form; encoding; codeword space}
  \T{Syndrome $s=He^T$: coset structure and syndrome decoding}
  \T{Minimum Hamming distance $d$; error detection and correction capability}
  \T{Repetition codes: $[n,1,n]$}
  \T{Hamming codes $[2^r-1, 2^r-1-r, 3]$: perfect codes, single-error correction}
  \T{Reed--Muller codes: recursive construction, puncturing, parameters}
  \T{Dual codes $C^\perp$; self-orthogonal and self-dual codes}
  \T{Singleton bound, Hamming bound, Plotkin bound, Gilbert--Varshamov bound}
  \T{LDPC codes (classical): Tanner graphs, parity checks of constant weight}
  \T{Belief propagation (BP) on Tanner graphs: message passing, sum-product algorithm}
  \T{Why dual-containing codes are central to CSS constructions}
}

\keybox{Construct and decode Hamming and Reed--Muller codes; run BP on a Tanner graph; state classical bounds and their quantum analogues.}

\newpage

%% ═══════════════════════════════════════════════════════════════
%%   PART II — ELEMENTARY QUANTUM CODES
%% ═══════════════════════════════════════════════════════════════

\parthead{II}{Elementary Quantum Codes}{8--14}
\addcontentsline{toc}{section}{Part II: Elementary Quantum Codes (Lec.\ 8--14)}

%%──────────────────────────────────────────────────────
\lecturetitle{8}{Three-Qubit Codes and the QEC Cycle}{First Codes}
\addcontentsline{toc}{subsection}{Lec.\ 8: Three-Qubit Codes and the QEC Cycle}

\topics{
  \T{Three-qubit bit-flip code $[[3,1,1]]$: encoding circuit $\ket{\psi}\to\ket{\psi\psi\psi}_{logical}$}
  \T{Syndrome extraction without measuring the logical qubit: ancilla-based parity checks}
  \T{Syndrome table and correction unitaries}
  \T{Three-qubit phase-flip code via Hadamard conjugation}
  \T{Fault paths and correlated errors (why two independent errors can fool the code)}
  \T{Resource overhead: 3 physical qubits per logical qubit, 2 ancillae}
  \T{Limitations: cannot correct both bit and phase flips simultaneously}
}

\keybox{Build and simulate the three-qubit QEC cycle in circuit form; identify uncorrectable error patterns.}

%%──────────────────────────────────────────────────────
\lecturetitle{9}{Shor's Nine-Qubit Code and Error Digitisation}{Concatenation}
\addcontentsline{toc}{subsection}{Lec.\ 9: Shor's Nine-Qubit Code and Error Digitisation}

\topics{
  \T{Concatenation strategy: combine bit-flip and phase-flip codes}
  \T{Encoding: $\ket{0_L}\to\tfrac{1}{2\sqrt{2}}(\ket{000}+\ket{111})^{\otimes 3}$, circuit derivation}
  \T{Syndrome extraction circuit: two rounds of parity measurements}
  \T{Correcting arbitrary single-qubit errors: $\{I,X,Y,Z\}$ form a basis---any error on one qubit correctable}
  \T{Error digitisation theorem: continuous errors project onto discrete Pauli errors upon syndrome measurement}
  \T{Resource count: 9 physical qubits, overhead analysis}
  \T{Connection to concatenated codes (preview of threshold theorem)}
}

\keybox{Derive the Shor code; prove error digitisation; count physical resources.}

%%──────────────────────────────────────────────────────
\lecturetitle{10}{The Knill--Laflamme Conditions}{QEC Theory}
\addcontentsline{toc}{subsection}{Lec.\ 10: The Knill--Laflamme Conditions}

\topics{
  \T{Formal definition of a quantum error-correcting code: subspace $\mathcal{C}\subseteq\mathcal{H}$}
  \T{Correctable error set $\mathcal{E}$: what it means}
  \T{Knill--Laflamme (KL) conditions: $\bra{c_i}E_a^\dagger E_b\ket{c_j}=\lambda_{ab}\delta_{ij}$ --- derivation}
  \T{Necessity and sufficiency of KL: proof of both directions}
  \T{Non-degenerate codes: $\lambda_{ab}=\lambda_a\delta_{ab}$; degenerate codes: $\lambda_{ab}$ non-diagonal}
  \T{Error spaces and orthogonality: geometric picture}
  \T{Syndrome measurement as the recovery map $\mathcal{R}$}
  \T{Quantum error correction as a projective recovery: Knill--Laflamme recovery channel}
  \T{Applications: verify KL for Shor code and three-qubit codes}
}

\keybox{State and prove the KL conditions in both directions; verify them for concrete codes; distinguish degenerate from non-degenerate codes.}

%%──────────────────────────────────────────────────────
\lecturetitle{11}{Quantum Bounds: Singleton, Hamming, and Beyond}{Bounds}
\addcontentsline{toc}{subsection}{Lec.\ 11: Quantum Bounds --- Singleton, Hamming, and Beyond}

\topics{
  \T{$[[n,k,d]]$ notation: $n$ physical qubits, $k$ logical qubits, distance $d$}
  \T{Quantum Singleton bound: $n - k \geq 4(d-1)$ --- derivation}
  \T{Quantum MDS codes: saturate the Singleton bound; existence and examples}
  \T{Quantum Hamming (sphere-packing) bound for non-degenerate codes: $\sum_{j=0}^{t}3^j\binom{n}{j}\leq 2^{n-k}$}
  \T{Perfect codes: saturate Hamming bound; uniqueness of $[[5,1,3]]$}
  \T{Quantum Plotkin bound}
  \T{Linear programming (LP) bounds: Delsarte/Krawtchouk framework extended to quantum codes}
  \T{Asymptotic rates: quantum Gilbert--Varshamov bound}
  \T{Degenerate codes and the open question of whether they surpass non-degenerate bounds}
}

\keybox{Derive quantum Singleton and Hamming bounds; assess a code's optimality; state the asymptotic rate-distance trade-off.}

%%──────────────────────────────────────────────────────
\lecturetitle{12}{The Five-Qubit Perfect Code $[[5,1,3]]$}{Perfect Code}
\addcontentsline{toc}{subsection}{Lec.\ 12: The Five-Qubit Perfect Code $[[5,1,3]]$}

\topics{
  \T{Existence and uniqueness: saturates the quantum Hamming bound}
  \T{Stabilizer generators: $XZZXI$, $IXZZX$, $XIXZZ$, $ZXIXZ$}
  \T{Logical operators $\bar{X}=XXXXX$, $\bar{Z}=ZZZZZ$}
  \T{Encoding circuit: systematic construction from stabilizers}
  \T{Syndrome table: all 15 non-trivial single-qubit errors}
  \T{Verification of KL conditions}
  \T{Decoding circuit}
  \T{Comparison with Shor code: same distance, lower overhead}
  \T{Transversal operations: which Clifford gates are transversal?}
}

\keybox{Construct the five-qubit code from scratch; verify KL; derive the complete syndrome table and encoding circuit.}

%%──────────────────────────────────────────────────────
\lecturetitle{13}{CSS Codes: The Calderbank--Shor--Steane Construction}{CSS}
\addcontentsline{toc}{subsection}{Lec.\ 13: CSS Codes --- The Calderbank--Shor--Steane Construction}

\topics{
  \T{Classical codes $C_1$ and $C_2$ with $C_2\subseteq C_1$; dual-containing condition $C_2^\perp \subseteq C_1$}
  \T{Construction: $X$-stabilizers from $H_2$, $Z$-stabilizers from $H_1$}
  \T{Parameters: $[[n, k_1 - k_2, d]]$ where $d=\min(d_1,d_2)$}
  \T{Independent $X$- and $Z$-error correction: decoupled syndrome decoding}
  \T{Transversal logical $\bar{X}$ and $\bar{Z}$; transversal CNOT}
  \T{Worked examples: CSS version of Shor code, Hamming-based CSS codes}
  \T{Importance for fault tolerance: separate $X/Z$ syndrome measurement}
}

\keybox{Construct any CSS code from two classical codes; identify stabilizers and parameters; explain why CSS enables transversal CNOTs.}

%%──────────────────────────────────────────────────────
\lecturetitle{14}{Quantum Codes from Algebraic Structures (Overview)}{Algebra}
\addcontentsline{toc}{subsection}{Lec.\ 14: Quantum Codes from Algebraic Structures (Overview)}

\topics{
  \T{Reed--Solomon codes and quantum Reed--Solomon (MDS family)}
  \T{BCH codes and quantum BCH codes}
  \T{Cyclic codes: generator polynomials and their quantum analogues}
  \T{Algebraic geometry (AG) codes: Goppa construction, quantum AG codes}
  \T{Quantum codes from combinatorial designs}
  \T{Convolutional quantum codes (brief)}
  \T{Quantum turbo codes (brief)}
  \T{Why these families matter: high rates, algebraic decoding algorithms}
}

\keybox{Describe how classical algebraic codes lift to quantum codes; state parameters of quantum BCH, RS, and AG codes.}

\newpage

%% ═══════════════════════════════════════════════════════════════
%%   PART III — STABILIZER FORMALISM
%% ═══════════════════════════════════════════════════════════════

\parthead{III}{Stabilizer Formalism}{15--22}
\addcontentsline{toc}{section}{Part III: Stabilizer Formalism (Lec.\ 15--22)}

%%──────────────────────────────────────────────────────
\lecturetitle{15}{The Pauli Group and Binary Symplectic Representation}{Stabilizers I}
\addcontentsline{toc}{subsection}{Lec.\ 15: The Pauli Group and Binary Symplectic Representation}

\topics{
  \T{Single-qubit Pauli group $\mathcal{P}_1 = \{\pm I, \pm iI, \pm X, \pm iX, \pm Y, \pm iY, \pm Z, \pm iZ\}$}
  \T{$n$-qubit Pauli group $\mathcal{P}_n$: structure, centre $\{\pm I, \pm iI\}^{\otimes n}$}
  \T{Commutation and anti-commutation relations; commutator $[P,Q]$ in $\mathcal{P}_n$}
  \T{Binary symplectic representation: $P \leftrightarrow (a|b)\in\mathbb{F}_2^{2n}$; $X\to(1|0)$, $Z\to(0|1)$, $Y\to(1|1)$}
  \T{Symplectic inner product $\Lambda(u,v) = u\cdot \Lambda \cdot v^T$: encodes commutativity}
  \T{Two Paulis commute iff symplectic inner product $=0$}
  \T{Check matrix $H=(H_X|H_Z)\in\mathbb{F}_2^{(n-k)\times 2n}$: rows are stabilizer generators}
  \T{Commutativity condition on $H$: $H_X H_Z^T + H_Z H_X^T = 0$ over $\mathbb{F}_2$}
}

\keybox{Represent any stabilizer group in binary symplectic form; test commutativity via the symplectic product.}

%%──────────────────────────────────────────────────────
\lecturetitle{16}{Stabilizer States and Stabilizer Codes}{Stabilizers II}
\addcontentsline{toc}{subsection}{Lec.\ 16: Stabilizer States and Stabilizer Codes}

\topics{
  \T{Stabilizer states: unique $+1$-eigenstate of an abelian stabilizer group}
  \T{Examples: computational basis states, Bell states, GHZ state, cluster states, graph states}
  \T{Stabilizer group of a code: $\mathcal{S}\subset\mathcal{P}_n$, abelian, $-I\notin\mathcal{S}$}
  \T{Code space $\mathcal{C} = \{\ket{\psi}: S\ket{\psi}=\ket{\psi}\;\forall S\in\mathcal{S}\}$; dimension $2^k$}
  \T{Logical operators: $\bar{X}_i, \bar{Z}_i$ commute with $\mathcal{S}$ but lie outside it}
  \T{Logical operator freedom: two sets of logical operators differ by elements of $\mathcal{S}$}
  \T{Parameters $[[n,k,d]]$: $k=n-m$ ($m$ generators), $d=$ min weight of logical operators}
  \T{Degenerate codes: stabilizer elements of weight $<d$}
  \T{Graph states and their stabilizer descriptions; the CZ-circuit construction}
}

\keybox{Identify the code space, logical operators, and parameters of any stabilizer code; construct graph states.}

%%──────────────────────────────────────────────────────
\lecturetitle{17}{Syndrome Measurement and Error Correction in Stabilizer Codes}{Stabilizers III}
\addcontentsline{toc}{subsection}{Lec.\ 17: Syndrome Measurement and Error Correction in Stabilizer Codes}

\topics{
  \T{Syndrome vector $s\in\mathbb{F}_2^m$: measuring each generator returns $\pm 1$}
  \T{Error $E$: syndrome $s_i = 0$ if $E$ commutes with $g_i$, $1$ if anti-commutes}
  \T{Syndrome measurement circuit: ancilla qubit per generator; CNOT/CZ fanout}
  \T{Why syndrome measurement does not disturb the logical qubit}
  \T{Coset structure: errors with the same syndrome form a coset of the normaliser}
  \T{Error correction: choose a representative from each coset (minimum-weight, or MLD)}
  \T{Residual logical error if wrong representative chosen}
  \T{Measurement errors: repeated rounds of syndrome measurement; space-time error model}
}

\keybox{Design syndrome extraction circuits for any stabilizer code; analyse measurement errors in repeated rounds.}

%%──────────────────────────────────────────────────────
\lecturetitle{18}{The Clifford Group and Gottesman Tableau Algorithm}{Clifford Group}
\addcontentsline{toc}{subsection}{Lec.\ 18: The Clifford Group and Gottesman Tableau Algorithm}

\topics{
  \T{Clifford group $\mathcal{C}_n$: unitaries mapping $\mathcal{P}_n\to\mathcal{P}_n$ under conjugation}
  \T{Generators: Hadamard $H$, phase gate $S$, CNOT; their actions on Paulis}
  \T{Clifford group modulo phases: symplectic group $\text{Sp}(2n,\mathbb{F}_2)$}
  \T{Gottesman--Knill theorem: Clifford circuits are classically efficiently simulatable}
  \T{Aaronson--Gottesman tableau: $2n \times 2n$ binary matrix tracking destabilizers and stabilizers}
  \T{Tableau update rules for $H$, $S$, CNOT, and Pauli measurements}
  \T{Complexity: $O(n^2)$ per gate, $O(n^3)$ for measurement}
  \T{Stim: a high-performance stabilizer circuit simulator (hands-on demo)}
  \T{Applications: simulating QEC circuits efficiently}
}

\keybox{Implement the Gottesman tableau; simulate circuits of hundreds of qubits classically; use Stim.}

%%──────────────────────────────────────────────────────
\lecturetitle{19}{Encoding and Decoding Circuits for Stabilizer Codes}{Circuits}
\addcontentsline{toc}{subsection}{Lec.\ 19: Encoding and Decoding Circuits for Stabilizer Codes}

\topics{
  \T{Standard form of the check matrix via Gaussian elimination over $\mathbb{F}_2$}
  \T{Systematic encoding circuit synthesis: algorithm and worked example}
  \T{Circuit depth and parallelism: scheduling syndrome measurements}
  \T{Decoding circuit: inverse of encoding unitary}
  \T{Ancilla qubit preparation: $\ket{0}$ vs.\ $\ket{+}$ ancillae for $Z$ and $X$ generators}
  \T{Circuit-level noise model: errors on gates, ancillae, and measurements}
  \T{Hands-on: implement the $[[5,1,3]]$ and $[[7,1,3]]$ encoding circuits in Stim}
}

\keybox{Systematically derive encoding and syndrome measurement circuits from any check matrix.}

%%──────────────────────────────────────────────────────
\lecturetitle{20}{The Steane $[[7,1,3]]$ Code: A Complete Study}{Steane Code}
\addcontentsline{toc}{subsection}{Lec.\ 20: The Steane $[[7,1,3]]$ Code --- A Complete Study}

\topics{
  \T{Hamming $[7,4,3]$ code: generator and parity-check matrices}
  \T{Dual-containing property: $[7,4,3]^\perp=[7,3,4]$; CSS construction}
  \T{Full stabilizer generators (6 generators): $H_Z$ and $H_X$ blocks}
  \T{Logical operators: $\bar{X}=X^7$, $\bar{Z}=Z^7$}
  \T{Encoding circuit (7 qubits, 6 ancillae)}
  \T{Syndrome extraction circuit and syndrome table (all 42 single-qubit error syndromes)}
  \T{Transversal Hadamard: $H^{\otimes 7}$ maps $\bar{X}\leftrightarrow\bar{Z}$}
  \T{Transversal CNOT between two code blocks}
  \T{Transversal phase gate $S^{\otimes 7}$ up to corrections}
  \T{Why Steane cannot transversally implement $T$: preview of Eastin--Knill}
}

\keybox{Fully construct and analyse the Steane code: stabilizers, logical operators, circuits, syndromes, and transversal gates.}

%%──────────────────────────────────────────────────────
\lecturetitle{21}{Syndrome Decoding Algorithms I: Classical and Minimum-Weight}{Decoding I}
\addcontentsline{toc}{subsection}{Lec.\ 21: Syndrome Decoding Algorithms I --- Classical and Minimum-Weight}

\topics{
  \T{Decoding problem: given syndrome $s$, find most likely error $\hat{E}$}
  \T{Lookup table (syndrome table) decoding: correctness, exponential storage}
  \T{Maximum likelihood decoding (MLD): $\hat{E}=\arg\max_E p(E|s)$ --- NP-hard in general}
  \T{Degenerate MLD: sum over degenerate cosets; coset-based decoding}
  \T{Minimum-weight (MW) decoding: polynomial-time for special codes}
  \T{Tanner graph of a stabilizer code: variable nodes (qubits), check nodes (generators)}
  \T{Belief propagation (BP) on the Tanner graph: log-likelihood ratio formulation}
  \T{BP failures: short cycles (girth), trapping sets}
  \T{Ordered statistics decoding (OSD): post-processing to fix BP failures}
  \T{BP-OSD algorithm and its performance for CSS codes}
  \T{Soft-decision decoding: exploiting analog measurement outcomes}
}

\keybox{Implement BP and BP-OSD decoders; analyse their complexity and failure modes; formulate MLD as an optimisation problem.}

%%──────────────────────────────────────────────────────
\lecturetitle{22}{Syndrome Decoding Algorithms II: Advanced Methods}{Decoding II}
\addcontentsline{toc}{subsection}{Lec.\ 22: Syndrome Decoding Algorithms II --- Advanced Methods}

\topics{
  \T{Small-set-flip (SSF) decoder: linear time, designed for qLDPC codes}
  \T{Renormalisation group (RG) decoder: multi-scale syndrome matching}
  \T{Maximum-likelihood decoding via tensor networks: exact MLD for 1D codes}
  \T{Neural network decoders: deep learning for decoding; training, architecture choices}
  \T{Reinforcement learning decoders (brief)}
  \T{Sliding-window decoding for real-time operation}
  \T{Space-time decoding: decoding in the presence of measurement errors}
  \T{Benchmarking decoders: threshold, sub-threshold logical error rate $p_L \propto (p/p_{th})^{\lceil d/2\rceil}$}
  \T{Decoder latency and real-time constraints for hardware}
}

\keybox{Compare advanced decoders on speed-accuracy trade-offs; implement and benchmark a neural-network decoder.}

\newpage

%% ═══════════════════════════════════════════════════════════════
%%   PART IV — FAULT-TOLERANT QUANTUM COMPUTING
%% ═══════════════════════════════════════════════════════════════

\parthead{IV}{Fault-Tolerant Quantum Computing}{23--30}
\addcontentsline{toc}{section}{Part IV: Fault-Tolerant Quantum Computing (Lec.\ 23--30)}

%%──────────────────────────────────────────────────────
\lecturetitle{23}{Principles of Fault Tolerance}{FT Principles}
\addcontentsline{toc}{subsection}{Lec.\ 23: Principles of Fault Tolerance}

\topics{
  \T{Error propagation through two-qubit gates: CNOT spreads $X$ forward, $Z$ backward}
  \T{Fault-tolerant (FT) gadget: a $t$-fault-tolerant gadget fails only if $\geq t+1$ faults occur}
  \T{Transversal gates: bitwise application to code blocks; why they cannot spread errors within a block}
  \T{Transversal $\neq$ FT in general: worked counter-examples}
  \T{Gadget composition: how FT gadgets are combined into FT circuits}
  \T{Malignant pairs: pairs of fault locations that together cause a logical error}
  \T{Level reduction: replacing a physical gadget with a higher-level FT gadget}
  \T{$t$-fault-tolerant protocols: $t$ faults $\Rightarrow$ at most $t$ errors on output block}
}

\keybox{Trace error propagation through any Clifford circuit; verify fault-tolerance of a given gadget; count malignant pairs.}

%%──────────────────────────────────────────────────────
\lecturetitle{24}{Fault-Tolerant Syndrome Extraction: Protocols and Trade-offs}{FT Syndrome}
\addcontentsline{toc}{subsection}{Lec.\ 24: Fault-Tolerant Syndrome Extraction --- Protocols and Trade-offs}

\topics{
  \T{Problem: a faulty ancilla can propagate errors into the data block}
  \T{Shor syndrome extraction: cat-state ancilla preparation + verification; analysis}
  \T{Steane syndrome extraction: encoded ancilla in the same code; transversal CNOT to data}
  \T{Knill (telecorrection) syndrome extraction: logical teleportation-based; removes data block from error path}
  \T{Flag qubit protocol: one or two flag qubits detect dangerous high-weight faults; resource efficient}
  \T{Comparison: ancilla overhead, depth, noise sensitivity, suitability for different hardware}
  \T{Repeated syndrome measurement: majority vote, minimum-weight recovery in space-time}
  \T{Circuit-level noise simulations: Monte Carlo with Stim}
}

\keybox{Describe, compare, and analyse Shor, Steane, Knill, and flag-qubit FT syndrome extraction; run circuit-level Monte Carlo simulations.}

%%──────────────────────────────────────────────────────
\lecturetitle{25}{Logical Gates I: Transversal and Clifford Gates}{Logical Gates I}
\addcontentsline{toc}{subsection}{Lec.\ 25: Logical Gates I --- Transversal and Clifford Gates}

\topics{
  \T{Transversal gates on CSS codes: $\bar{H}=H^{\otimes n}$, $\bar{S}=S^{\otimes n}$ (where applicable), $\overline{CNOT}=CNOT^{\otimes n}$}
  \T{Which codes support which transversal Clifford generators}
  \T{Gate teleportation: implementing a logical gate via teleportation with an ancilla state $U\ket{0}$}
  \T{Teleportation-based Clifford gates: cost and ancilla preparation}
  \T{Code switching (gauge fixing): changing between two codes to access different transversal gate sets}
  \T{Example: switching between Steane and Reed--Muller codes to access $\bar{T}$ and $\bar{H}$}
  \T{Pieceable fault tolerance: gates that are not transversal but can be done in FT pieces}
}

\keybox{Implement transversal Clifford gates; explain gate teleportation; describe and execute code switching.}

%%──────────────────────────────────────────────────────
\lecturetitle{26}{Logical Gates II: The Eastin--Knill Theorem}{Eastin-Knill}
\addcontentsline{toc}{subsection}{Lec.\ 26: Logical Gates II --- The Eastin--Knill Theorem}

\topics{
  \T{Statement: no quantum error-correcting code with continuous symmetry has a transversal universal gate set}
  \T{Precise statement: no QECC admits a universal transversal gate set}
  \T{Proof sketch: transversal gates form a group; a finite group cannot be universal}
  \T{Implications: must use non-transversal methods to achieve universality}
  \T{Approximate versions and exceptions: codes with approximate transversal $T$ gates}
  \T{Circumventing Eastin--Knill: magic state injection, code switching, gauge fixing, 3D codes}
  \T{Bravyi--König theorem: 2D topological codes have transversal gates in the third level of the Clifford hierarchy at most}
}

\keybox{Prove the Eastin--Knill theorem; identify strategies to bypass it; state the Bravyi--König theorem.}

%%──────────────────────────────────────────────────────
\lecturetitle{27}{Magic State Distillation}{Magic States}
\addcontentsline{toc}{subsection}{Lec.\ 27: Magic State Distillation}

\topics{
  \T{Magic states: $\ket{T}=T\ket{+}$ and $\ket{H}=H\ket{0}$; why they suffice for universality with Clifford gates}
  \T{Noisy magic states: $\rho_T = (1-p)\proj{T} + p \frac{I}{2}$; threshold for distillation}
  \T{Bravyi--Kitaev (BK) 15-to-1 distillation protocol: uses $[[15,1,3]]$ Reed--Muller code; derivation}
  \T{Error analysis: output fidelity as a function of input error rate $p$}
  \T{Distillation threshold: $p < p_{th}\approx 0.14$ for BK protocol}
  \T{Resource cost: 15 noisy $\ket{T}$ per output; $k$ levels of concatenation give $p^{2^k}$}
  \T{Improved protocols: 10-to-2 (Meier--Eastin--Knill), punctured Reed--Muller, triorthogonal codes}
  \T{Triorthogonal codes: codes whose generators satisfy a triorthogonality condition; enable transversal $T$}
  \T{Resource estimation: physical qubit and time costs for distillation factories}
  \T{Magic state distillation factories in surface code architecture}
}

\keybox{Derive the BK distillation protocol; compute output error rate; estimate factory overhead for a target error rate.}

%%──────────────────────────────────────────────────────
\lecturetitle{28}{Concatenated Codes and the Threshold Theorem}{Threshold}
\addcontentsline{toc}{subsection}{Lec.\ 28: Concatenated Codes and the Threshold Theorem}

\topics{
  \T{Concatenated codes: encoding a logical qubit with code $C$, then encoding each physical qubit with $C$ again}
  \T{Parameters of level-$\ell$ concatenation: $[[n^\ell, 1, d^\ell]]$}
  \T{Logical error rate under concatenation: $p_L^{(\ell)} \approx \frac{1}{c}\left(c\,p\right)^{2^\ell}$ for distance-3 codes}
  \T{Threshold error rate $p_{th} = 1/c$: below threshold, $p_L\to 0$ with $\ell$}
  \T{Rigorous threshold theorem: Aharonov--Ben-Or; Kitaev; Knill--Laflamme--Zurek}
  \T{Proof sketch (Aharonov--Ben-Or): level reduction, malignant pair counting}
  \T{Resource overhead: $O(\text{poly}\log(1/\epsilon))$ physical operations per logical gate at error $\epsilon$}
  \T{Physical vs.\ pseudo-threshold: code-specific one-level breakeven point}
  \T{Numerical thresholds for common codes: $\sim10^{-4}$ (depolarizing, concatenated), $\sim1\%$ (surface codes)}
  \T{Limitations: correlated noise, non-Markovian errors, leakage, threshold reduction}
}

\keybox{Prove the threshold theorem via level reduction; compute resource overhead; discuss limitations.}

%%──────────────────────────────────────────────────────
\lecturetitle{29}{Fault-Tolerant Protocol Design and Resource Estimation}{FT Design}
\addcontentsline{toc}{subsection}{Lec.\ 29: Fault-Tolerant Protocol Design and Resource Estimation}

\topics{
  \T{Full FT architecture: logical qubit storage, syndrome extraction, gate application, magic state factories}
  \T{Space-time cost of a logical circuit: $T$-count, $T$-depth, CNOT-count}
  \T{Resource estimation frameworks: (e.g., Beverland et al., Azure Quantum Resource Estimator)}
  \T{Estimating physical qubits for Shor's algorithm: $\sim 10^6$ to $10^8$ physical qubits at $10^{-3}$ error rate}
  \T{Algorithmic-level optimisations: $T$-gate synthesis, Clifford+$T$ decomposition}
  \T{Rotation synthesis: Solovay--Kitaev algorithm; $\epsilon$-approximation of arbitrary single-qubit gates}
  \T{Gridsynth and Repeat-Until-Success (RUS) circuits for efficient rotation synthesis}
  \T{Co-design considerations: choosing code and hardware jointly}
}

\keybox{Perform end-to-end resource estimation for a quantum algorithm; apply Solovay--Kitaev and gridsynth for rotation synthesis.}

%%──────────────────────────────────────────────────────
\lecturetitle{30}{Quantum Error Correction Experiments: State of Practice}{Experiments}
\addcontentsline{toc}{subsection}{Lec.\ 30: Quantum Error Correction Experiments --- State of Practice}

\topics{
  \T{Superconducting qubit experiments: Google ($[[5,1,3]]$, surface code below threshold, 2023--24), IBM}
  \T{Trapped-ion experiments: IonQ, Quantinuum ($[[7,1,3]]$ Steane, color codes, logical qubit lifetime)}
  \T{Neutral atom experiments: Harvard/QuEra, Pasqal}
  \T{Photonic approaches: GKP states in optical modes}
  \T{Metrics: logical error rate per round, break-even point, logical qubit lifetime vs.\ physical $T_1$}
  \T{Current achievements and open engineering challenges}
  \T{Roadmap to fault-tolerant quantum computing: timeline estimates}
}

\keybox{Compare QEC achievements across platforms; assess where each platform stands relative to fault-tolerance thresholds.}

\newpage

%% ═══════════════════════════════════════════════════════════════
%%   PART V — TOPOLOGICAL AND SURFACE CODES
%% ═══════════════════════════════════════════════════════════════

\parthead{V}{Topological and Surface Codes}{31--37}
\addcontentsline{toc}{section}{Part V: Topological and Surface Codes (Lec.\ 31--37)}

%%──────────────────────────────────────────────────────
\lecturetitle{31}{Topological Quantum Codes: Motivation and the Toric Code}{Toric Code}
\addcontentsline{toc}{subsection}{Lec.\ 31: Topological Quantum Codes --- Motivation and the Toric Code}

\topics{
  \T{Why topology? Local stabilizers + high threshold + 2D hardware compatibility}
  \T{Kitaev's toric code: qubits on edges of a square lattice on a torus}
  \T{Vertex operators $A_v = \prod_{e\ni v}X_e$ and plaquette operators $B_p = \prod_{e\in p}Z_e$}
  \T{Ground state space: $[[2L^2, 2, L]]$ code; two logical qubits from torus topology}
  \T{Logical operators: non-contractible loops (strings) in $X$ and $Z$ bases}
  \T{Anyons: $e$ (electric, $A_v=-1$) and $m$ (magnetic, $B_p=-1$) excitations}
  \T{Anyon creation, propagation, and annihilation as string operators}
  \T{Braiding statistics: $e$ and $m$ anyons are mutual semions}
  \T{Topological protection: logical error $\Leftrightarrow$ system-spanning string}
}

\keybox{Construct the toric code; identify anyons and their braiding statistics; relate topological order to logical operators.}

%%──────────────────────────────────────────────────────
\lecturetitle{32}{The Planar (Surface) Code: Construction and Properties}{Surface Code I}
\addcontentsline{toc}{subsection}{Lec.\ 32: The Planar (Surface) Code --- Construction and Properties}

\topics{
  \T{Open boundary conditions: rough ($Z$) and smooth ($X$) boundaries}
  \T{Stabilizers on a $d\times d$ grid: $\lfloor d^2/2\rfloor$ each of $A_v$ and $B_p$; one logical qubit}
  \T{Logical operators: horizontal string ($\bar{X}$), vertical string ($\bar{Z}$); weight $d$}
  \T{Code parameters: $[[d^2+(d-1)^2, 1, d]]$ (rotated) or $[[2d^2-2d+1, 1, d]]$ (unrotated)}
  \T{Rotated planar code: checkerboard layout; $d^2$ data qubits, $(d^2-1)$ ancillae}
  \T{Syndrome measurement circuit: weight-4 and weight-2 boundary stabilizers}
  \T{Circuit-level noise model; $X$ and $Z$ errors and measurement errors}
  \T{Logical error rate vs.\ $d$ and $p$: exponential suppression below threshold}
}

\keybox{Construct and parameterise the rotated planar code; draw syndrome extraction circuits; describe the error model.}

%%──────────────────────────────────────────────────────
\lecturetitle{33}{Surface Code Threshold and Logical Error Rate}{Surface Code II}
\addcontentsline{toc}{subsection}{Lec.\ 33: Surface Code Threshold and Logical Error Rate}

\topics{
  \T{Heuristic threshold argument: matching random errors to random syndrome changes}
  \T{Mapping to 2D random-bond Ising model (Nishimori line)}
  \T{Code-capacity threshold $\approx 10.9\%$ (depolarizing); circuit-level threshold $\approx 0.5$--$1\%$}
  \T{Sub-threshold logical error rate: $p_L \approx A(p/p_{th})^{\lfloor d/2 \rfloor + 1}$}
  \T{Numerical Monte Carlo estimation with Stim + PyMatching}
  \T{Teraquop qubit: physical qubits needed for one logical qubit at $10^{-6}$ logical error rate}
  \T{Biased noise: tailored surface codes and XZZX surface codes}
}

\keybox{Derive and interpret the surface code threshold; simulate logical error rates numerically; understand tailored-code advantages.}

%%──────────────────────────────────────────────────────
\lecturetitle{34}{Surface Code Decoding: MWPM and Union-Find}{Decoding III}
\addcontentsline{toc}{subsection}{Lec.\ 34: Surface Code Decoding --- MWPM and Union-Find}

\topics{
  \T{Decoding problem: matching syndrome defects to cancel them in a minimum-weight way}
  \T{Syndrome graph: nodes are defects (anyons), edge weights from error probabilities}
  \T{Minimum-weight perfect matching (MWPM): formulation as a graph problem}
  \T{Blossom V algorithm: $O(n^3)$ exact MWPM; implementation outline}
  \T{Weighted matching with log-likelihood weights under independent noise}
  \T{Space-time matching: including measurement errors; 3D syndrome graph}
  \T{Union-Find (Helios/Fowler) decoder: near-linear time $O(n\alpha(n))$; algorithm}
  \T{Accuracy vs.\ speed: MWPM vs.\ Union-Find trade-off}
  \T{Correlated matching decoders: accounting for $Y$ errors (correlation between $X$ and $Z$ decoding)}
  \T{PyMatching: open-source MWPM library (hands-on demo)}
}

\keybox{Implement MWPM and Union-Find decoders; benchmark them on Monte Carlo simulations; compare correlated and uncorrelated decoding.}

%%──────────────────────────────────────────────────────
\lecturetitle{35}{Surface Code Decoding: Neural Networks and Beyond}{Decoding IV}
\addcontentsline{toc}{subsection}{Lec.\ 35: Surface Code Decoding --- Neural Networks and Beyond}

\topics{
  \T{Motivation: MWPM is sub-optimal under non-Pauli or correlated noise}
  \T{Convolutional neural network (CNN) decoders: training on syndrome patterns}
  \T{Recurrent and transformer-based decoders}
  \T{Reinforcement learning (RL) decoders: agent learns to correct syndromes}
  \T{Tensor network decoders: exact MLD for 2D codes via contraction (approximate)}
  \T{Renormalisation group (RG) decoder: $O(n\log n)$ multi-scale matching}
  \T{Accuracy comparison: neural decoders can approach MLD; at cost of training data and latency}
  \T{Real-time decoding requirements: decoder must keep pace with syndrome generation ($\mu$s regime)}
  \T{FPGA and ASIC implementations of decoders}
}

\keybox{Train and evaluate a neural network decoder; compare decoder accuracy and latency; discuss real-time implementation.}

%%──────────────────────────────────────────────────────
\lecturetitle{36}{Logical Gates on Surface Codes: Lattice Surgery}{Lattice Surgery}
\addcontentsline{toc}{subsection}{Lec.\ 36: Logical Gates on Surface Codes --- Lattice Surgery}

\topics{
  \T{Transversal gates on surface codes: only $\bar{X}$ and $\bar{Z}$ are transversal}
  \T{Lattice surgery: merge and split operations on adjacent code patches}
  \T{Logical CNOT via lattice surgery: merge two patches, measure a joint stabilizer, split}
  \T{Logical CNOT circuit: ancilla patch, measurement sequence, correction}
  \T{Logical $\bar{H}$ via patch rotation (code deformation) or twist defects}
  \T{Logical $\bar{S}$ via lattice surgery with a $\ket{Y}$ ancilla state}
  \T{Logical $\bar{T}$ via magic state injection: distilled $\ket{T}$ state + CNOT + measurement + correction}
  \T{Space-time cost of lattice surgery operations}
  \T{3D space-time diagrams: ZX-calculus representation of lattice surgery}
}

\keybox{Design lattice surgery protocols for all Clifford gates; estimate space-time costs; draw ZX-calculus diagrams.}

%%──────────────────────────────────────────────────────
\lecturetitle{37}{Colour Codes, Honeycomb Codes, and Other 2D Topological Codes}{Other 2D Codes}
\addcontentsline{toc}{subsection}{Lec.\ 37: Colour Codes, Honeycomb Codes, and Other 2D Topological Codes}

\topics{
  \T{Colour codes on three-colourable lattices: 4.8.8 and 6.6.6 lattices}
  \T{Colour code stabilizers: $X$ and $Z$ on each face; parameters $[[n, k, d]]$}
  \T{Transversal gate set on colour codes: full Clifford group transversally; $\bar{T}$ transversally on 3D colour codes}
  \T{Folding map: colour code is equivalent to two stacked surface codes}
  \T{Decoding colour codes: restriction decoder, Möbius decoder}
  \T{Honeycomb (Floquet) code (Hastings--Haah): dynamical code driven by two-qubit Pauli measurements; no static stabilizers}
  \T{Measurement schedule of the honeycomb code: six-step cycle generating effective topological order}
  \T{Hastings--Haah--O'Brien code: low overhead for superconducting hardware with limited connectivity}
  \T{Comparison: 2D code families on threshold, transversal gate set, overhead, and hardware compatibility}
}

\keybox{Describe colour code structure and its richer transversal gate set; explain Floquet codes; compare all 2D topological code families.}

\newpage

%% ═══════════════════════════════════════════════════════════════
%%   PART VI — QUANTUM LDPC CODES
%% ═══════════════════════════════════════════════════════════════

\parthead{VI}{Quantum LDPC Codes}{38--41}
\addcontentsline{toc}{section}{Part VI: Quantum LDPC Codes (Lec.\ 38--41)}

%%──────────────────────────────────────────────────────
\lecturetitle{38}{Classical LDPC Codes and the Tanner Graph: Deep Dive}{Classical LDPC}
\addcontentsline{toc}{subsection}{Lec.\ 38: Classical LDPC Codes and the Tanner Graph --- Deep Dive}

\topics{
  \T{LDPC code definition: parity-check matrix with constant row and column weight}
  \T{Tanner graph: bipartite graph; variable nodes (bits) and check nodes (parities)}
  \T{Girth and its effect on BP performance: longer cycles $\Rightarrow$ better BP}
  \T{Expander codes: spectral expansion and its role in distance}
  \T{Sipser--Spielman linear-time decoding for expander codes}
  \T{Capacity-approaching LDPC codes: irregular degree distributions, density evolution}
  \T{Spatially coupled LDPC codes: threshold saturation}
  \T{Why LDPC matters for quantum: constant-overhead fault tolerance possible}
}

\keybox{Analyse LDPC codes via their Tanner graphs; understand expansion and its role in decoding; explain capacity-approaching properties.}

%%──────────────────────────────────────────────────────
\lecturetitle{39}{Quantum LDPC Codes: Hypergraph Product and Beyond}{qLDPC I}
\addcontentsline{toc}{subsection}{Lec.\ 39: Quantum LDPC Codes --- Hypergraph Product and Beyond}

\topics{
  \T{From CSS to LDPC: commutativity constraint; challenge of keeping both check matrices sparse}
  \T{Hypergraph product (HGP) code (Tillich--Z{\'e}mor): $H = [H_1\otimes I, I\otimes H_2^T]$ construction}
  \T{Parameters of HGP codes: $[[n_1 n_2 + n_1' n_2', k_1 k_2 + k_1' k_2', \min(d_1,d_2)]]$; constant rate if input codes have constant rate}
  \T{Lifted product (LP) codes: generalization using quasi-cyclic structure}
  \T{Fiber bundle codes (Hastings--Haah--O'Brien): twisting the fiber to boost distance}
  \T{Balanced product codes: another generalization}
  \T{Distance vs.\ rate: current landscape; $[[n, \Theta(n^{1/2}), \Theta(n^{1/2})]]$ achievable}
}

\keybox{Construct HGP and LP codes; compute their parameters; describe the code parameter landscape.}

%%──────────────────────────────────────────────────────
\lecturetitle{40}{Good Quantum LDPC Codes: Breakthroughs and Constructions}{qLDPC II}
\addcontentsline{toc}{subsection}{Lec.\ 40: Good Quantum LDPC Codes --- Breakthroughs and Constructions}

\topics{
  \T{Definition of ``good'' codes: $[[n, \Theta(n), \Theta(n)]]$ with constant-weight checks}
  \T{Quantum expander codes (Leverrier--Tillich--Z{\'e}mor): $[[n, \Theta(n), \Theta(\sqrt{n})]]$}
  \T{Panteleev--Kalachev (2022): first asymptotically good qLDPC codes; lifted product over chain complexes}
  \T{Leverrier--Z{\'e}mor quantum Tanner codes (2022): construction, parameters, and proof sketch}
  \T{Chain complex perspective: codes as 2-complexes; boundary operators $\partial_1, \partial_2$}
  \T{Cohomological products: cup product and product of complexes}
  \T{Hyperbolic codes and codes on higher-dimensional manifolds}
  \T{Consequences for overhead: constant-overhead fault tolerance theorem (Gottesman, 2014; Leverrier et al.)}
}

\keybox{Describe the Panteleev--Kalachev and Leverrier--Z{\'e}mor constructions; explain the chain complex framework; state the constant-overhead fault tolerance theorem.}

%%──────────────────────────────────────────────────────
\lecturetitle{41}{Decoding Quantum LDPC Codes and Hardware Considerations}{qLDPC III}
\addcontentsline{toc}{subsection}{Lec.\ 41: Decoding Quantum LDPC Codes and Hardware Considerations}

\topics{
  \T{Decoding qLDPC codes: BP on the quantum Tanner graph; failure for short cycles}
  \T{BP-OSD for qLDPC: performance and complexity}
  \T{Small-set-flip (SSF) decoder (Leverrier--Z{\'e}mor): linear-time; correctness proof sketch}
  \T{Sliding-window decoding for real-time qLDPC}
  \T{Neural decoders for qLDPC codes}
  \T{Hardware realisation of qLDPC: non-local connectivity requirement; long-range interactions}
  \T{Reconfigurable neutral atom arrays: shuttling enables non-local operations}
  \T{Interconnected modules (modular architectures): quantum buses, teleportation links}
  \T{Overhead comparison: qLDPC vs.\ surface codes; $O(1)$ overhead vs.\ $O(d^2)$ overhead}
  \T{Current state: smallest experimental qLDPC demonstrations}
}

\keybox{Implement and benchmark SSF and BP-OSD for a qLDPC code; compare overhead with surface codes; discuss hardware feasibility.}

\newpage

%% ═══════════════════════════════════════════════════════════════
%%   PART VII — BOSONIC CODES AND HARDWARE
%% ═══════════════════════════════════════════════════════════════

\parthead{VII}{Bosonic Codes and Hardware}{42--44}
\addcontentsline{toc}{section}{Part VII: Bosonic Codes and Hardware (Lec.\ 42--44)}

%%──────────────────────────────────────────────────────
\lecturetitle{42}{Bosonic Codes: Cat Qubits and Binomial Codes}{Bosonic I}
\addcontentsline{toc}{subsection}{Lec.\ 42: Bosonic Codes --- Cat Qubits and Binomial Codes}

\topics{
  \T{Bosonic encoding: logical qubit in a harmonic oscillator (cavity mode)}
  \T{Fock space: $\{|0\rangle, |1\rangle, |2\rangle, \ldots\}$; creation/annihilation operators $\hat{a}^\dagger, \hat{a}$}
  \T{Noise model for oscillators: photon loss (amplitude damping), dephasing, photon gain}
  \T{Coherent states and cat states: $|\pm\alpha\rangle$; two-component cat code}
  \T{Kerr-cat qubit: Kerr Hamiltonian confines system to two-cat-state manifold; exponential $Z$-bias}
  \T{Dissipative cat qubit: engineered two-photon dissipation stabilises cat states}
  \T{Bias-preserving gates on cat qubits: CNOT and $Z$-rotation without dephasing}
  \T{Binomial codes: $|0_L\rangle$ and $|1_L\rangle$ as superpositions of Fock states; designed for photon loss}
  \T{Pair-cat codes, four-component cat codes}
  \T{Concatenation of cat qubits with qubit codes: biased-noise repetition code; exponential suppression}
}

\keybox{Describe cat qubit stabilisation; explain bias-preserving gates; design a concatenated cat+repetition code.}

%%──────────────────────────────────────────────────────
\lecturetitle{43}{GKP Codes and Grid State Encodings}{Bosonic II}
\addcontentsline{toc}{subsection}{Lec.\ 43: GKP Codes and Grid State Encodings}

\topics{
  \T{Gottesman--Kitaev--Preskill (GKP) code: encoding a qubit in a harmonic oscillator via modular variables}
  \T{Stabilizers: displacement operators $e^{i2\sqrt{\pi}\hat{q}}$ and $e^{-i2\sqrt{\pi}\hat{p}}$}
  \T{Logical operators: $\bar{X}=e^{i\sqrt{\pi}\hat{q}}$, $\bar{Z}=e^{-i\sqrt{\pi}\hat{p}}$}
  \T{Ideal GKP state: unphysical comb of delta functions; approximate GKP states}
  \T{Error correction: homodyne measurement of $\hat{q}$ mod $\sqrt{\pi}$; displacement correction}
  \T{GKP Steane-type correction: teleportation-based correction}
  \T{Noise model: Gaussian displacement errors; effective qubit error rate}
  \T{GKP code performance: near-hashing bound for Gaussian noise}
  \T{Hardware: superconducting microwave cavities (cQED), optical modes, trapped-ion motion}
  \T{Rectangular and hexagonal GKP lattices; multi-mode GKP codes}
  \T{Concatenation of GKP with qubit codes: GKP-surface code, GKP-repetition code}
}

\keybox{Derive GKP stabilizers and logical operators; describe error correction via modular homodyne measurement; design concatenated GKP+qubit codes.}

%%──────────────────────────────────────────────────────
\lecturetitle{44}{Hardware Platforms and Co-design of Codes}{Hardware \& Co-design}
\addcontentsline{toc}{subsection}{Lec.\ 44: Hardware Platforms and Co-design of Codes}

\topics{
  \T{Superconducting qubits (transmons): gate times, connectivity (nearest-neighbour grid), $T_1/T_2$, leakage}
  \T{Superconducting cavities (3D cQED): photon lifetime $>1$\,ms; GKP and cat qubit host}
  \T{Trapped ions: all-to-all connectivity, long coherence times, slow gates; suited for Steane/colour codes}
  \T{Neutral atoms in tweezers: reconfigurable connectivity (shuttling), mid-circuit measurement; suited for qLDPC}
  \T{Photonic quantum computing: cluster-state approach, linear optics, GKP in optical modes}
  \T{Co-design principles: choose code to match hardware noise bias, connectivity graph, and gate set}
  \T{Noise bias exploitation: XZZX surface code for biased noise; cat+repetition for pure $Z$-bias}
  \T{Connectivity-aware syndrome circuits: reducing non-local operations}
  \T{Compiling logical circuits to hardware: routing, scheduling, and decomposition}
  \T{System-level resource estimation: from algorithm to physical qubits}
}

\keybox{Match code families to hardware platforms; perform co-design given a noise model; compile logical circuits to hardware-native gates.}

\newpage

%% ═══════════════════════════════════════════════════════════════
%%   PART VIII — RESEARCH FRONTIERS
%% ═══════════════════════════════════════════════════════════════

\parthead{VIII}{Research Frontiers and Open Problems}{45}
\addcontentsline{toc}{section}{Part VIII: Research Frontiers and Open Problems (Lec.\ 45)}

%%──────────────────────────────────────────────────────
\lecturetitle{45}{Open Problems and the Road Ahead}{Frontiers}
\addcontentsline{toc}{subsection}{Lec.\ 45: Open Problems and the Road Ahead}

\topics{
  \T{\textbf{Decoding frontiers:} optimal decoders for qLDPC codes; efficient exact MLD; decoding under coherent noise}
  \T{\textbf{New code constructions:} single-shot error correction codes (no repeated measurements); high-rate codes for specific hardware; quantum polar codes}
  \T{\textbf{Single-shot QEC:} Bombin's 3D codes; codes where one round of noisy measurements suffices for correction}
  \T{\textbf{Quantum fault tolerance with non-Pauli noise:} coherent errors, non-Markovian noise, crosstalk; approximate QEC}
  \T{\textbf{Approximate quantum error correction (AQEC):} trading exact correctability for better parameters; Petz recovery map; quantum error mitigation}
  \T{\textbf{Quantum error mitigation (QEM):} zero-noise extrapolation, probabilistic error cancellation, Clifford data regression; relation to and distinction from QEC}
  \T{\textbf{Quantum memory:} self-correcting quantum memories; 4D toric code; Haah's cubic code; open problem: does a self-correcting quantum memory exist in 3D?}
  \T{\textbf{Constant-overhead fault tolerance:} practical implementations of qLDPC-based architectures}
  \T{\textbf{Quantum LDPC codes for entanglement distillation and quantum communications}}
  \T{\textbf{Connections to condensed matter:} topological phases of matter, anyons, SPT phases, quantum spin liquids}
  \T{\textbf{Connections to quantum gravity and holography:} quantum error correction in AdS/CFT; Ryu--Takayanagi formula as QEC; tensor network codes}
  \T{\textbf{Classical simulation of QEC:} limits of stabilizer simulation; beyond-Clifford simulation}
  \T{How to navigate the research literature; key conferences (QIP, TQC, APS); important research groups}
}

\keybox{Identify and articulate open problems in QEC; connect QEC to condensed matter and quantum gravity; navigate the research literature independently.}

\newpage

%% ═══════════════════════════════════════════════════════════════
%%   APPENDIX: COURSE ADMINISTRATION
%% ═══════════════════════════════════════════════════════════════

\section*{Course Administration}
\addcontentsline{toc}{section}{Course Administration}

\subsection*{Grading}
\begin{center}
\begin{tabular}{llc}
  \toprule
  \textbf{Component} & \textbf{Description} & \textbf{Weight} \\
  \midrule
  Problem Sets (8 total)     & Biweekly; proofs, circuit design, simulations & 40\% \\
  Midterm Exam               & Covers Parts I--III; closed book                & 15\% \\
  Final Exam                 & Comprehensive; closed book                      & 20\% \\
  Programming Projects (2)   & Decoder implementation; resource estimation     & 15\% \\
  Research Presentation      & Paper from last 3 years; 20-min talk + Q\&A     & 10\% \\
  \bottomrule
\end{tabular}
\end{center}

\subsection*{Software Tools Used in the Course}
\begin{itemize}[itemsep=3pt]
  \item \textbf{Stim} (Google) --- high-performance stabilizer circuit simulator; used from Lecture 18 onward
  \item \textbf{PyMatching} --- MWPM decoder for surface codes; used in Lectures 34--35
  \item \textbf{LDPC} (Python) --- BP-OSD for classical and quantum LDPC codes; Lectures 21, 41
  \item \textbf{Qiskit QEC} --- IBM's error correction toolkit; supplementary exercises
  \item \textbf{NumPy / SciPy / Matplotlib} --- numerical experiments and visualisation throughout
  \item \textbf{CVXPY} --- SDP for diamond norm computations; Lecture 5
\end{itemize}

\subsection*{Primary References}

\noindent\textbf{Textbooks:}
\begin{itemize}[itemsep=2pt]
  \item Nielsen \& Chuang, \textit{Quantum Computation and Quantum Information} (Cambridge UP, 2000) --- Ch.\ 10, 12
  \item Lidar \& Brun (eds.), \textit{Quantum Error Correction} (Cambridge UP, 2013)
  \item Gottesman, \textit{Surviving as a Quantum Computer in a Classical World} (lecture notes, 2024 edition)
  \item Preskill, \textit{Lecture Notes on Quantum Computation}, Chapter 7 (online)
\end{itemize}

\noindent\textbf{Key Research Papers (a representative sample):}
\begin{itemize}[itemsep=2pt, label={$\triangleright$}]
  \item Shor (1995); Steane (1996) --- founding codes
  \item Knill \& Laflamme (1997) --- KL conditions
  \item Gottesman (1997) --- stabilizer formalism (thesis)
  \item Kitaev (2003) --- fault-tolerant quantum computation by anyons
  \item Dennis, Kitaev, Landahl, Preskill (2002) --- topological quantum memory
  \item Fowler, Martinis et al.\ (2012) --- surface codes: towards practical large-scale QC
  \item Bravyi \& Kitaev (2005) --- magic state distillation
  \item Eastin \& Knill (2009) --- restrictions on transversal gates
  \item Hastings \& Haah (2021) --- honeycomb code
  \item Panteleev \& Kalachev (2022) --- asymptotically good quantum LDPC codes
  \item Leverrier \& Z{\'e}mor (2022) --- quantum Tanner codes
  \item Gottesman (2014) --- fault tolerance with constant overhead
  \item Grimsmo \& Puri (2021) --- GKP codes (review)
  \item Google Quantum AI (2023) --- below-threshold surface code experiment
\end{itemize}

\subsection*{Lecture--Topic Quick Reference}

{\small
\begin{longtable}{@{}cll c@{}}
  \toprule
  \textbf{Lec.} & \textbf{Title} & \textbf{Key Concept} & \textbf{Part} \\
  \midrule
  \endfirsthead
  \toprule
  \textbf{Lec.} & \textbf{Title} & \textbf{Key Concept} & \textbf{Part} \\
  \midrule
  \endhead
  \bottomrule
  \endfoot
  1  & Case for QEC                    & No-cloning; error digitisation          & I \\
  2  & Density Operators               & Mixed states, partial trace             & I \\
  3  & CPTP Maps \& Kraus Operators    & Choi iso., Stinespring dilation         & I \\
  4  & Noise Channels                  & Bit-flip, dephasing, depolarizing, AD   & I \\
  5  & Distance Measures               & Fidelity, trace dist., diamond norm     & I \\
  6  & Hardware Characterisation       & RB, QPT, GST                           & I \\
  7  & Classical Coding Theory         & BP on Tanner graphs, dual codes         & I \\
  \midrule
  8  & Three-Qubit Codes               & Syndrome without disturbing logical     & II \\
  9  & Shor's Code                     & Error digitisation; concatenation       & II \\
  10 & Knill--Laflamme Conditions      & Necessity and sufficiency proof         & II \\
  11 & Quantum Bounds                  & Singleton, Hamming, LP bounds           & II \\
  12 & Five-Qubit Perfect Code         & Unique $[[5,1,3]]$; KL verification     & II \\
  13 & CSS Construction                & Dual-containing codes; transversal CNOT & II \\
  14 & Algebraic Code Families         & BCH, RS, AG quantum codes               & II \\
  \midrule
  15 & Pauli Group \& Symplectic Rep.  & Binary symplectic inner product         & III \\
  16 & Stabilizer States \& Codes      & Code space, logical operators           & III \\
  17 & Syndrome Extraction             & Ancilla-based non-destructive readout   & III \\
  18 & Clifford Group \& Tableaux      & Gottesman--Knill; Aaronson--Gottesman   & III \\
  19 & Encoding \& Decoding Circuits   & Systematic circuit synthesis            & III \\
  20 & Steane $[[7,1,3]]$ Code         & Full worked study                       & III \\
  21 & Decoding Algorithms I           & BP, BP-OSD, MLD                         & III \\
  22 & Decoding Algorithms II          & NN decoders, SSF, sliding window        & III \\
  \midrule
  23 & FT Principles                   & Error propagation; malignant pairs      & IV \\
  24 & FT Syndrome Extraction          & Shor, Steane, Knill, flag qubits        & IV \\
  25 & Logical Gates I: Clifford       & Gate teleportation; code switching      & IV \\
  26 & Eastin--Knill Theorem           & Proof; Bravyi--König theorem            & IV \\
  27 & Magic State Distillation        & BK 15-to-1; triorthogonal codes         & IV \\
  28 & Threshold Theorem               & Concatenated codes; level reduction     & IV \\
  29 & FT Protocol Design              & Resource estimation; gridsynth          & IV \\
  30 & QEC Experiments                 & Hardware milestones; open challenges    & IV \\
  \midrule
  31 & Toric Code                      & Anyons; braiding statistics             & V \\
  32 & Planar Surface Code             & Rotated code; syndrome circuits         & V \\
  33 & Surface Code Threshold          & Ising mapping; Monte Carlo; $p_L(d,p)$  & V \\
  34 & Decoding: MWPM \& Union-Find    & Blossom V; $O(n\alpha(n))$ UF decoder   & V \\
  35 & Decoding: Neural \& Beyond      & CNN, RL, tensor network decoders        & V \\
  36 & Lattice Surgery                 & Logical gates via merge/split           & V \\
  37 & Colour \& Honeycomb Codes       & Floquet codes; 2D code comparison       & V \\
  \midrule
  38 & Classical LDPC Deep Dive        & Expander codes; density evolution       & VI \\
  39 & Hypergraph Product Codes        & HGP, LP, fiber bundle constructions     & VI \\
  40 & Good qLDPC Codes                & Panteleev--Kalachev; Leverrier--Z{\'e}mor & VI \\
  41 & qLDPC Decoding \& Hardware      & SSF; neutral atom realisation           & VI \\
  \midrule
  42 & Cat Qubits \& Binomial Codes    & Kerr-cat; dissipative cat; bias-preserving & VII \\
  43 & GKP Codes                       & Grid states; homodyne correction        & VII \\
  44 & Hardware Platforms \& Co-design & Platform comparison; noise-code matching & VII \\
  \midrule
  45 & Open Problems \& Research Frontiers & Single-shot; AQEC; self-correcting memories & VIII \\
\end{longtable}
}

\vfill
\begin{center}
  \color{rulegray}\small
  \rule{0.5\textwidth}{0.4pt}\\[4pt]
  \textit{This outline reflects the state of the field as of 2025--26.
  Parts VI--VIII track active research; content is updated annually.}
\end{center}

\end{document}
